\section{Formal Ingredients}
\label{sec:model}

The paper uses two semantic axes and one update discipline.

\begin{itemize}
  \item Capability answers \emph{where} the current computation runs.
  \item Region answers \emph{where} a non-\code{Sendable} value is currently connected.
  \item Affine environment update answers \emph{what} remains usable after evaluation.
\end{itemize}

The key claim is that Swift Concurrency becomes predictable only when these three roles are kept distinct.

\subsection{Judgment and Declaration Boundaries}

\begin{rulebox}{Expression judgment}
\[
  \Gamma;\; @\kappa;\; \alpha \vdash e : T \at \rho \dashv \Gamma'
\]

\begin{itemize}
  \item $\Gamma$: input environment
  \item $\Gamma'$: output environment
  \item $@\kappa$: current execution capability
  \item $\alpha$: ambient mode, either $\mathit{sync}$ or $\mathit{async}$
  \item $\rho$: result region
\end{itemize}

The final component is essential because evaluation may preserve a binding, refine its region, or remove it from the caller-side environment.
\end{rulebox}

Swift function declarations are not expression terms, so the formalization also needs a declaration judgment that says how a body is checked. A declaration contributes two pieces of information to its body: the body capability~$@\kappa$ and the ambient mode~$\alpha$. For a function type annotation $@\varphi = @\sigma\; @\iota$, the body capability is obtained from $\toCapability(\mathrm{proj}_\iota(@\varphi))$, while $\alpha$ is determined by the presence or absence of \code{async}. This is where the paper introduces \code{await}-permission; it is not encoded as a separate effect system.

\begin{definitionbox}{Function annotations and body capability}
\[
\begin{aligned}
  @\sigma &::= \cdot \mid @\textsf{Sendable} \\
  @\iota &::= @\kappa \mid @\textsf{isolated(any)} \mid @\textsf{concurrent} \\
  @\varphi &::= @\sigma\; @\iota
\end{aligned}
\]

\[
\begin{aligned}
  \toCapability(@\kappa) &= @\kappa \\
  \toCapability(@\textsf{isolated(any)}) &= @\textsf{nonisolated} \\
  \toCapability(@\textsf{concurrent}) &= @\textsf{nonisolated}
\end{aligned}
\]
\end{definitionbox}

This split between $@\sigma$ and $@\iota$ matters immediately. \code{@Sendable} is not an execution mode. It constrains capture. By contrast, $@\iota$ determines the body capability. Both \code{@concurrent async} and \code{nonisolated(nonsending) async} therefore type-check their bodies under $@\textsf{nonisolated}$, even though their call-site semantics differ later~\cite{se0461}.

\subsection{Capability, Region, and Two Access Relations}

Capability alone does not explain what remains legal after a call or an assignment. The model also needs a value-side location discipline. The region component is inspired by Tofte and Talpin's region-based memory management~\cite{tofte1997region}, while the capability component follows the static-capability tradition of Walker, Crary, and Morrisett~\cite{walker2000capabilities}. The interplay between regions and linear ownership draws on Walker and Watkins~\cite{walker2001regions}.

\[
\begin{aligned}
  \rho_{\mathit{ns}} &::= \disconnected \mid \isolated(a) \mid \task \mid \invalid \\
  \rho &::= \rho_{\mathit{ns}} \mid \mathunderscore
\end{aligned}
\]

$\disconnected$ means not yet bound to an isolation domain. $\isolated(a)$ means connected to a concrete actor region. $\task$ denotes task-local non-\code{Sendable} state. $\mathunderscore$ is the normalized region tag for \code{Sendable} values. The normal form condition is therefore simple: $T : \textsf{Sendable}$ if and only if $\rho = \mathunderscore$.

Two relations are then needed.

\begin{itemize}
  \item $\Accessible(@\kappa)$ determines which regions may be used directly under capability $@\kappa$.
  \item $\Capturable(@\kappa)$ determines which regions may be closed over by a closure checked under $@\kappa$.
\end{itemize}

This distinction is one of the paper's critical refinements. Actor-isolated code may directly access task-bound values in some contexts, but an actor-isolated closure may not capture a \code{task} value because the captured state could outlive the dynamic situation that made the direct access safe. That is why nonisolated async parameters can be captured by nonisolated and \code{@isolated(any)} closures yet still be rejected by \code{@MainActor} closures after SE-0461~\cite{se0461}.

\subsection{Region Merge and Capability-Region Glue}

\codefile{snippets/region-binding.swift}

The snippet above illustrates why region and capability cannot be merged into one notion. The value \code{x} begins disconnected. Storing it into actor-isolated state does not change its nominal type, but it does change the future facts that the type checker must remember. The environment must therefore record a refinement from $\disconnected$ to an actor-bound region.

\paperfigure[0.78]{figures/region-merge-semilattice.png}{Region merge for non-\code{Sendable} values. A disconnected value can be refined into an actor or task region, while incompatible branches join to $\invalid$, which corresponds to a compile-time failure.}

The operation that connects capability and region is $\toRegion(@\kappa)$. It answers a narrow but important question: if a same-isolation operation binds a value to the current execution context, which region should the value acquire? For a concrete actor capability the answer is that actor's isolated region; for $@\textsf{nonisolated}$ the answer remains $\disconnected$. This definition is what lets the formalization state same-isolation binding without conflating execution and storage.

\begin{propositionbox}{Why the output environment is unavoidable}
Region merge is not a side remark. It is the mechanism that explains why local code changes future legality. The result of an assignment or a same-isolation call is often \code{Sendable} and therefore normalized to region $\mathunderscore$, while the interesting fact survives only in the updated environment. A model without $\Gamma'$ can describe the value just produced, but it cannot describe why a later transfer becomes illegal.
\end{propositionbox}

\subsection{Transfer Is a Predicate, Not a Slogan}

Swift offers both explicit and implicit transfer. Explicit transfer is written with \code{sending}. Implicit transfer occurs when a non-\code{Sendable} value crosses to a distinct isolation domain. The paper packages both cases into one predicate:

\[
  \canSend(@\kappa, @\iota, [\sending])
\]

$\canSend$ holds either because the parameter is explicitly marked \code{sending}, or because the callee isolation $@\iota$ is distinct from the caller capability $@\kappa$ and therefore induces an implicit cross-isolation transfer. This is another place where surface syntax hides a deeper commonality: the important fact is not merely whether the keyword \code{sending} appears, but whether the call site transfers responsibility for the value.

This viewpoint also clarifies the asymmetry between parameters and results. Parameters can be consumed from the caller by shrinking the environment. Results travel in the opposite direction. A cross-isolation non-\code{Sendable} result is therefore legal only if the callee explicitly promises a disconnected \code{sending} result. The next two sections show how these auxiliary definitions govern conversion, calls, closure inference, and dynamic isolation.
